% Recovered verbatim from chapters_90_95_with_part/chapters_90_95_with_part.tex
% Canonical chapter 95
\chapter{Adaptive Phase Interfaces and Assistive Cognitive Technologies}

\section{Adaptive Interface Field}
An assistive interface couples to the cognitive field via
\[
\tau_{\mathrm{int}}(x,t) = \alpha\,\tau_{\mathrm{cog}}(x,t) + \beta\,\tau_{\mathrm{env}}(x,t),
\]
and adjusts its output to minimise
\[
\mathcal{E}_{\mathrm{int}} = \int (\partial\tau_{\mathrm{int}} - \partial\tau_{\mathrm{target}})^2 dV,
\]
where \(\tau_{\mathrm{target}}\) encodes a desired comfort or focus state.

\section{Toward Neuro-Supportive Devices}
By continuously estimating \(\tau_{\mathrm{cog}}\) and \(\tau_{\mathrm{env}}\), adaptive phase interfaces can in principle:
\begin{itemize}
  \item damp disruptive gradients (e.g.\ sensory overload),
  \item enhance stabilising modes (e.g.\ focus, calm),
  \item and provide a real-time, geometry-based feedback channel between mind and environment.
\end{itemize}
This chapter lays the groundwork for future MK--43 Ultra inspired assistive technologies, including advanced auditory and multisensory modulation systems.
